\subsection{2. Xử lý tập dữ liệu và Cài đặt tiền xử lý}

\subsubsection{2.1. Tập dữ liệu Minh-TFR và Phân bố mẫu}
\indent\indent Tập dữ liệu thực nghiệm được xây dựng từ tập chuẩn TFR và DTD-lightweight, bao gồm các cặp ảnh tài liệu định dạng PNG/JPG, mặt nạ nhị phân ground-truth kích thước tương ứng và các nhãn thuộc tính làm giả được lưu trữ dưới định dạng JSON Lines (\texttt{.jsonl}):
\begin{itemize}[label={--}, itemsep=2pt]
    \item \textbf{Tập dữ liệu huấn luyện (\texttt{train\_datas.jsonl}):} Chứa hàng ngàn mẫu tài liệu đa dạng như biên lai, hóa đơn điện tử, hợp đồng, chứng minh thư.
    \item \textbf{Tập dữ liệu kiểm thử (\texttt{test\_datas.jsonl}):} Tập kiểm thử độc lập \textbf{Minh-TFR TestSet} gồm 500 batches đánh giá chuẩn dùng để đo lường năng lực tổng quát hóa.
\end{itemize}

\subsubsection{2.2. Thuật toán phân tích nhãn và Lọc cụm từ phủ định}
\indent\indent Trong các tập dữ liệu chú thích văn bản, nhiều mẫu xuất hiện các câu mô tả ngữ cảnh phủ định (ví dụ: \textit{``no obvious edge artifacts'', ``not cause inconsistent font'', ``should be deleted''}). Để tránh gán nhầm nhãn dương tính ($1$), thuật toán trích xuất vector nhãn sử dụng danh sách các cụm từ loại trừ:
\begin{lstlisting}[language=Python, frame=single, caption={Thuật toán trích xuất vector nhãn 6 chiều}]
def extract_label_vector(content):
    vec = []
    for k in LABELS:
        m = re.search(LABEL_PATTERNS[k], content, re.I)
        if m:
            text = m.group(1).lower()
            if not any(neg in text for neg in NEGATIVE_PHRASES):
                vec.append(1)
            else:
                vec.append(0)
        else:
            vec.append(0)
    return vec
\end{lstlisting}

\subsubsection{2.3. Pipeline tăng cường dữ liệu đồng bộ ảnh -- mặt nạ}
\indent\indent Do đặc thù bài toán phân đoạn pixel, các phép biến đổi hình học phải được áp dụng đồng bộ cho cả ảnh đầu vào và mặt nạ ground-truth:
\begin{enumerate}
    \item \textbf{Biến đổi hình học đồng thời (Coupled Geometric Transforms):} Áp dụng phép lật ngang ngẫu nhiên (\texttt{RandomHorizontalFlip}, $p=0.5$) và xoay nhẹ (\texttt{RandomRotation} trong khoảng $\pm 10^\circ$) đồng thời cho cả ảnh và mặt nạ.
    \item \textbf{Biến đổi quang học (Photometric Transforms):} Áp dụng \texttt{ColorJitter} (độ sáng, độ tương phản $\pm 0.1$) chỉ riêng cho ảnh đầu vào, giữ nguyên mặt nạ nhị phân.
    \item \textbf{Chuẩn hóa ImageNet:} Chuẩn hóa kênh màu ảnh theo $\mu = [0.485, 0.456, 0.406]$ và $\sigma = [0.229, 0.224, 0.225]$. Lưu ý rằng khi đưa vào mô-đun BlockDCT, ảnh được chuyển đổi ngược về dạng xám trong dải $[0, 1]$ trước khi tính toán phổ tần số để đảm bảo tính ổn định số học.
\end{enumerate}
